\section{总结与延伸}

\subsection{从 Behavior 到 Mechanism 的完整链条}

本讲最值得复用的不是某张 attention heatmap，而是一套研究程序：先定义行为指标，再寻找 abrupt change；选择可解 toy model 建立 algebra；从 path structure 提出候选 algorithm；最后用 ablation、timing、cross-example generality 与反例给证据分级。Mechanistic interpretability 的可信度来自这条链条，而不是一张“看起来像在关注正确 token”的图。

\begin{importantbox}{十五条核心结论}
\begin{enumerate}
  \item Mechanistic interpretability 尝试把 parameters 反编译为 algorithms。
  \item ICL 固定参数，只在 context 与 activations 中发生临时适应。
  \item 二维 loss surface 同时展示 training snapshot 与 token position。
  \item $\mathrm{Loss@500}-\mathrm{Loss@50}$ 越负，late-context advantage 越强。
  \item $0.4$ nats 约为 $0.58$ bits/token，不是 40\% accuracy。
  \item Loss bump 与 derivative spike 产生 phase-change hypothesis，但不构成证明。
  \item Attention-only toy models 去掉 MLP/LayerNorm/bias 以换取精确 algebra。
  \item Head 可分为 QK routing 与 OV writing 两个 circuit。
  \item 固定 attention pattern 后，attention-only path 对输入局部线性。
  \item One-layer models 主要实现 skip-trigram/positional rules，ICL 弱且无 phase change。
  \item Eigenvalue lens 可总结 small token-to-token maps，但在 MLP/大模型中失效。
  \item Two-layer expansion 产生 composed/virtual heads，并允许 routing 读取前层 features。
  \item Induction circuit 实现 $[A][B]\ldots[A]\to[B]$ 的 context nearest neighbor。
  \item Small-model ablation 是因果证据；large-model formation timing 主要是相关证据。
  \item Induction 是重要机制而非必然唯一机制，soft induction 与 distributed circuits 仍需更强验证。
\end{enumerate}
\end{importantbox}

\subsection{一个可复用的 Circuit Audit Protocol}

面对新的可解释性主张，可以依次检查：目标 behavior 是否量化；候选 component 是否有结构能力；路径是否跨 layers 组合；attention pattern 与 OV logit effect 是否一致；ablation 是否保留 distribution 与其他 pattern；结果是否跨 prompts/seeds/models；结论是否区分 toy model、small model 与 production-scale model。

\begin{knowledgebox}{实践检查表}
\begin{enumerate}
  \item 定义 behavior metric 与 sign convention；
  \item 画 training × context 的二维诊断，而非只看最终 checkpoint；
  \item 展开 residual paths，区分 routing 与 writing；
  \item 用 activation/attention 找 hypothesis，用 ablation 做因果检查；
  \item 记录 compensation、redundancy 与 negative-result heads；
  \item 对 large-model extrapolation 标注 correlation、plausibility 或 intervention；
  \item 主动寻找 alternative mechanism 和 failure case。
\end{enumerate}
\end{knowledgebox}

\subsection{拓展阅读}

建议先读 \emph{A Mathematical Framework for Transformer Circuits} 的 one-/two-layer decomposition，再读 \emph{In-context Learning and Induction Heads} 的 evidence taxonomy 与 appendix。随后对照 Kaplan scaling laws 和 GPT-3 few-shot behavior，思考 behavioral curve、training dynamics 与 circuit formation 的关系。阅读现代大型模型解释工作时，应重点观察它们如何处理 MLP nonlinearities、feature superposition、normalization、distributed computation 与 causal validation，而不是只比较 head visualizations 是否漂亮。

\teachervoice{本讲最后的 present theory 是两个主要 algorithms：广义 copying 与 induction-head in-context nearest neighbors，同时承认可能还有别的机制。把这份不确定性保留下来，正是对课堂与原始研究最忠实的总结。}

\end{document}
